% !TeX program = xelatex
% !TeX encoding = UTF-8
% RSFI — статья ВАК К2 в стандарте группы (шаблон v8: xelatex, Times New Roman).
% Сборка: xelatex -interaction=nonstopmode rsfi_vak.tex (дважды).
\documentclass[10pt,a4paper]{article}
\usepackage[T1]{fontenc}
\usepackage{fontspec}
\setmainfont{Times New Roman}
\usepackage[russian]{babel}
\usepackage{amsmath,amssymb,amsthm}
\usepackage{booktabs}
\usepackage{graphicx}
\setlength{\abovecaptionskip}{2pt}
\setlength{\belowcaptionskip}{1pt}
\usepackage[a4paper,margin=1.7cm,top=1.65cm,bottom=1.65cm]{geometry}
\linespread{0.95}
\newcommand{\secfmt}[1]{\vspace{5pt}\noindent{\large\bfseries #1}\par\vspace{1.5pt}}
\newcommand{\subsecfmt}[1]{\vspace{4pt}\noindent{\normalsize\bfseries #1}\par\vspace{1.5pt}}
\usepackage{url}
\urlstyle{tt}
\setlength{\parskip}{1.0pt}
\setlength{\parindent}{1.25em}
\setlength{\tabcolsep}{3.5pt}

\newtheorem{proposition}{Предложение}
\newtheorem{theorem}{Теорема}
\newtheorem{corollary}{Следствие}
\newtheorem{lemma}{Лемма}
\newtheorem{definition}{Определение}

\begin{document}
\noindent{\small\bfseries УДК 004.056:004.852}\par\vspace{4pt}
\begin{center}
{\large\bfseries Safe-aware дискриминантная коррекция и внутриклассовое отбеливание для быстрых геометрических фильтров детекции jailbreak-атак (семейство RSFI)\par}
\vspace{4pt}
Осанов~В.А.$^{1}$, Лебедев~П.Д.$^{1}$, Ворфоломеев~С.В.$^{1}$\\
{\small $^{1}$Поволжский государственный университет телекоммуникаций и информатики (ПГУТИ), 443010, Россия, г.~Самара, Московское шоссе, д.~77\\
e-mail: \texttt{me@a3ldi.ru}}
\end{center}
\vspace{3pt}

\noindent\textbf{Аннотация.} Работа посвящена сверхбыстрым геометрическим фильтрам первого эшелона (L1) для детекции jailbreak-атак на большие языковые модели в пространстве текстовых эмбеддингов. Показано, что одноклассовые фильтры (наивный косинус, SVD-подпространство) систематически теряют качество на гетерогенных данных, поскольку игнорируют геометрию безопасного класса. Предложены: (1)~safe-aware дискриминантная коррекция $B1$ — направление $\mu_{mal}-\mu_{safe}$ (один вектор, $O(d)$ на запрос), дающая прирост +7.7--10.5~п.п.\ ROC-AUC над наивным косинусом на гетерогенном корпусе WildChat ($p<0{,}0001$, парный тест ДеЛонга, 5~сидов, воспроизводится на четырёх эмбеддерах размерности 768--4096); (2)~внутриклассовое отбеливание $\Sigma_W^{-1/2}$, для которого доказано точное разложение $\Sigma_T=\Sigma_W+\Sigma_B$ с межклассовой добавкой ранга~1 и формула деградации тотального отбеливания через тождество Шермана--Моррисона; $\Sigma_W$-отбеливание восстанавливает в среднем 43{,}9\% потерь тотального отбеливания и не уступает ему ни в одной из 12~конфигураций. Метод $B1$ превосходит опубликованные трансформерные классификаторы ProtectAI DeBERTa-v3 (0{,}8405) и unitary/toxic-bert (0{,}7236) на WildChat при задержке скоринга около 0{,}003~мс против единиц--десятков миллисекунд. Честно зафиксированы границы применимости: статическая обфускация (base64) и адаптивный жадный подбор синонимов (обход при изменении менее 3\% слов, коразмерность~1); рекомендована роль ультрабыстрого эшелона L1 в архитектуре defense-in-depth, а не замена нейросетевым судьям. Все числа воспроизводимы из закоммиченных CSV и покрыты 248~автотестами.\par\vspace{2pt}
\noindent\textbf{Ключевые слова:} безопасность больших языковых моделей; jailbreak; prompt injection; эмбеддинги; геометрические фильтры; дискриминантное направление; отбеливание; оценка Ледуа--Вольфа; разложение ковариации; guardrails.\par\vspace{5pt}

\begin{center}
{\normalsize\bfseries Safe-Aware Discriminant Correction and Pooled Within-Class Whitening for Fast Geometric Filters in LLM Jailbreak Detection (RSFI Family)\par}
\vspace{2pt}
{\small Osanov~V.A.$^{1}$, Lebedev~P.D.$^{1}$, Vorfolomeev~S.V.$^{1}$\\
$^{1}$Povolzhskiy State University of Telecommunications and Informatics (PSUTI), 77 Moskovskoe sh., Samara, 443010, Russia\\
e-mail: \texttt{me@a3ldi.ru}}
\end{center}
\vspace{2pt}

\noindent\textbf{Abstract.} We study ultrafast first-echelon (L1) geometric filters for jailbreak-attack detection on large language models in text-embedding space. One-class filters (naive cosine, SVD subspace) systematically lose accuracy on heterogeneous data because they ignore the geometry of the safe class. We propose: (1)~safe-aware discriminant correction $B1$ — the direction $\mu_{mal}-\mu_{safe}$ (a single vector, $O(d)$ per query), gaining +7.7--10.5~pp ROC-AUC over naive cosine on the heterogeneous WildChat corpus ($p<0{,}0001$, paired DeLong test, 5~seeds, reproduced on four embedders of dimension 768--4096); (2)~pooled within-class whitening $\Sigma_W^{-1/2}$, for which we prove the exact decomposition $\Sigma_T=\Sigma_W+\Sigma_B$ with a rank-1 between-class term and a total-whitening degradation formula via the Sherman--Morrison identity; $\Sigma_W$-whitening recovers on average 43{,}9\% of total-whitening losses and never loses to it in any of 12~configurations. Method $B1$ outperforms published transformer classifiers ProtectAI DeBERTa-v3 (0{,}8405) and unitary/toxic-bert (0{,}7236) on WildChat at a scoring latency of about 0{,}003~ms versus units to tens of milliseconds. Applicability boundaries are honestly reported: static obfuscation (base64) and adaptive greedy synonym substitution (bypass at under 3\% word changes, codimension~1); we recommend the role of an ultrafast L1 echelon in a defense-in-depth architecture rather than a replacement for neural judges. All numbers are reproducible from committed CSV files and covered by 248~automated tests.\par\vspace{2pt}
\noindent\textbf{Keywords:} large language model safety; jailbreak; prompt injection; embeddings; geometric filters; discriminant direction; whitening; Ledoit--Wolf estimation; covariance decomposition; guardrails.\par\vspace{4pt}
\hrule\vspace{5pt}
\secfmt{1. Введение}

Внедрение больших языковых моделей (LLM) в критические корпоративные контуры (финтех, здравоохранение, юридический сектор) упирается в уязвимости к состязательным атакам: jailbreak, инъекции промптов, семантические манипуляции [11, 26]. Зрелые enterprise-решения к 2026~году сошлись к трёхслойной архитектуре защиты («Three-Clocks»): L1~--- быстрые геометрические/статистические фильтры эмбеддингов (доли миллисекунды), L2~--- специализированные классификаторы (десятки--сотни миллисекунд), L3~--- LLM-судьи для асинхронного аудита (секунды). Настоящая работа занимает уровень~L1: фильтр, который отсекает очевидные атаки до того, как запрос достигнет дорогих эшелонов L2/L3.

Проблема известных L1-фильтров двояка. Во-первых, одноклассовые методы (косинус до центроида атак, SVD-подпространство угроз) слепы к безопасному классу и деградируют на стилистически гетерогенных данных. Во-вторых, сферическое отбеливание, устраняющее анизотропию эмбеддингов («семантический конус» [3, 2]), помогает не всегда: на гетерогенных корпусах оно вредит разделению, и до настоящей работы причина этого не была формализована.

\textit{О названии и рамках.} Историческое имя проекта~--- Riemannian System Fidelity Index (RSFI). Риманов аппарат (логарифмическая и экспоненциальная карты на $\mathbb{S}^{d-1}$) реализован в библиотеке (\path|src/rsfi/geometry.py|, класс \path|RiemannianSphere|) и покрыт юнит-тестами, однако бенчмаркированный в работе конвейер использует евклидову сферическую геометрию: $L_2$-нормализацию, ZCA-отбеливание Ледуа--Вольфа и SVD/дискриминантные направления. Это разграничение строгое и осознанное; никаких римановых карт в оценённых методах нет.

\textbf{Вклад работы:}

\begin{enumerate}\setlength{\itemsep}{0pt}\setlength{\parsep}{0pt}
\item Safe-aware дискриминантная коррекция $B1$ (направление $\mu_{mal}-\mu_{safe}$, один вектор, $O(d)$): +7.7--10.5~п.п.\ ROC-AUC на WildChat, 5/5~сидов, $p<0{,}0001$; уровень, близкий к supervised-потолку (логистическая регрессия).
\item Теория внутриклассового отбеливания: точное тождество $\Sigma_T=\Sigma_W+\Sigma_B$, формула подавления сигнала через Шермана--Моррисона, метод $B1w$ ($\Sigma_W^{-1/2}$) со средним восстановлением 43{,}9\% потерь.
\item Прямое сравнение с опубликованными классификаторами на общих сплитах: лёгкий фильтр превосходит DeBERTa-v3 и toxic-bert при задержке на три порядка ниже.
\item Честные границы: обфускация (base64: $B1$ сохраняет 0{,}705 против 0{,}238 у наивного косинуса), адаптивный противник (коразмерность~1, обход при $<3\%$ замен), кросс-доменный перенос и операционные точки (TPR при FPR~1\%).
\end{enumerate}

\secfmt{2. Обзор литературы}

\textbf{Геометрия представлений.} Анизотропия контекстных эмбеддингов (концентрация в узком конусе, ненулевое среднее) установлена в [3]; отбеливание представлений предложений~--- в [2]. Риманова оптимизация на многообразиях изложена в [4, 5, 6]. Направление отказа как геометрический объект («refusal is mediated by a single direction») показано в [7]; раздельное кодирование вредоносности и отказа~--- в [8]. Стриминговая детекция по траекториям скрытых состояний (TrajGuard) [9] и инженерия представлений [10] развивают смежный white-box фланг, тогда как настоящая работа остаётся строго black-box (только эмбеддинги запросов).

\textbf{Защита LLM.} Базовые guardrail-подходы: Llama Guard [11], NeMo Guardrails [12], позиционная статья о guardrails [13], модерационные API. Лёгкие эмбеддинговые защиты (drift detection, steering, codebooks) активно исследуются [14--17]. Отличие настоящей работы: строгий leakage-free протокол с фиксированными сидами, тестами ДеЛонга и автотестами консистентности каждой ячейки отчёта, плюс доказанная теория отбеливания и честно измеренные границы, включая адаптивного осведомлённого противника (в духе [24]).

\textbf{Данные и метрики.} Корпусы: ToxicChat [18] (5082~реальных диалога: 384~вредоносных, 4698~безопасных), WildChat (гетерогенный открытый корпус; рабочая выборка 2000~запросов: 1000/1000), XSTest [19] (450~контрастных омонимических пар: 200~unsafe / 250~safe), AdvBench [20] (520~вредоносных) и HarmBench [21] (400) в паре с 400~безопасными Alpaca [22]. Сравнение ROC-кривых~--- парный тест ДеЛонга [23]; оценка ковариаций при $N_{ref}\le d$~--- усадка Ледуа--Вольфа [1].

\secfmt{3. Математическая модель}

\subsecfmt{3.1. Латентное пространство как сфера}

Пусть $\mathcal{X}$~--- множество последовательностей токенов над конечным алфавитом, $E:\mathcal{X}\to\mathbb{R}^d$~--- энкодер ($d\in\{768,1024,4096\}$, $d\ge 2$). Нормализация
\[
\hat{x}=\frac{E(x)}{\|E(x)\|_2}
\]
отображает состояния на единичную гиперсферу
\[
\mathcal{M}=\mathbb{S}^{d-1}=\{x\in\mathbb{R}^d:\|x\|_2=1\},
\]
с геодезическим расстоянием (длиной дуги большой окружности)
\[
d_{\mathcal{M}}(x,y)=\arccos\langle x,y\rangle\in[0,\pi].
\]
Эмбеддинги естественного языка анизотропны: $\mu=\mathbb{E}[E(x)]\ne\mathbf{0}$ (концентрация в узком конусе). По калибровочному корпусу $\{x_i\}_{i=1}^N$ оцениваются среднее $\hat{\mu}$ и ковариация $\Sigma$; при $N_{ref}\le d$ вырожденность снимается усадкой Ледуа--Вольфа $\Sigma_{LW}=(1-\rho)\Sigma+\rho\mu I$, $\mu=\tfrac1d\mathrm{Tr}(\Sigma)$, $\rho\in[0,1]$ оптимален. ZCA-отбеливание $W=\Sigma_{LW}^{-1/2}=V\Lambda^{-1/2}V^T$ с последующей ренормализацией
\[
\tilde{x}=\frac{W(x-\mu)}{\|W(x-\mu)\|_2}\in\mathbb{S}^{d-1}
\]
выравнивает масштабы дисперсии перед скорингом.

\subsecfmt{3.2. Дискриминантное направление}

Одноклассовый SVD-базис на вредоносных векторах находит направления \textit{внутренней} дисперсии класса угроз, а не разделяющую гиперплоскость. Двухклассовая коррекция использует центроиды обоих классов:
\[
\mathbf{w}_{disc}=\frac{\mu_{mal}-\mu_{safe}}{\|\mu_{mal}-\mu_{safe}\|_2},\qquad
s_{B1}(x)=\tilde{x}^T\mathbf{w}_{disc},
\]
одно скалярное произведение, $O(d)$, один хранимый вектор. Контрастивный SVD $(X_{mal}-\mu_{safe})^T=U_{cont}S_{cont}V_{cont}^T$~--- родственное $k$-мерное обобщение; RSFI-функционал касательного пространства ($\|\mathbf{v}_\perp\|-\alpha\pi_{thr}-\beta d_{\mathcal{M}}$, решающее правило $\mathrm{BLOCK}$ при $\mathrm{rsfi}<\tau$) в продакшен-конвенции знака ранг-эквивалентен $B1$ ($|\Delta\mathrm{AUC}|\le 0{,}0068$ на всех 9~конфигурациях; в линейном пределе $\alpha\to\infty$, $\beta=0$ эквивалентность точная с отклонением $\le 6{,}1\cdot10^{-3}$).

\subsecfmt{3.3. Разложение ковариации и семейство Махаланобис-дискриминантов}

\begin{definition}[Сферический ZCA-отбеливатель]
Для SPD-матрицы $\Sigma=V\Lambda V^T$: $W(\Sigma)=\Sigma^{-1/2}$, $T_\Sigma(x)=W(\Sigma)(x-\mu_\Sigma)$, $\hat{x}=T_\Sigma(x)/\|T_\Sigma(x)\|$.
\end{definition}

\begin{proposition}[ZCA-коммутация отбеливания и дискриминанта]
Направление дискриминанта по отбелённым пулам совпадает с отбелённым дискриминантом исходного пространства: $T_\Sigma(\mu_m)-T_\Sigma(\mu_s)=W(\Sigma)d$, $d=\mu_m-\mu_s$, независимо от весов смеси. Итоговый скор~--- сферический Махаланобис-косинус $B1_\Sigma(x)=\mathrm{smc}_{\Sigma^{-1}}(x;d)$; семейство $\{B1,B1b,B1w\}$~--- один метод с метриками $I$, $\Sigma_T^{-1}$, $\Sigma_W^{-1}$.
\end{proposition}

\begin{theorem}[Точное тождество]
Для смеси двух классов с весами $(\pi_m,\pi_s)$:
\[
\boxed{\Sigma_T=\Sigma_W+\Sigma_B,\qquad \Sigma_B=\pi_m\pi_s\,dd^T,}
\]
тотальная ковариация равна пулу внутриклассовых плюс ранг-1 межклассовая добавка вдоль $d$ (закон полной ковариации).
\end{theorem}

\begin{corollary}[Формула деградации]
По Шерману--Моррисону:
\[
\boxed{\|d\|_{\Sigma_T^{-1}}=\frac{\|d\|_{\Sigma_W^{-1}}}{\sqrt{1+\pi_m\pi_s\|d\|_{\Sigma_W^{-1}}^2}}.}
\]
Отбеливание по $\Sigma_T$ подавляет именно используемый $B1$ сигнал; потеря растёт с внутриклассовым Махаланобис-расстоянием $r_W=\|d\|_{\Sigma_W^{-1}}$: при $r_W\gg1$ (гетерогенный режим) тотальное отбеливание разрушительно, при $r_W\lesssim1$ (анизотропный фон, омонимия) безвредно.
\end{corollary}

\begin{lemma}[Валидность при усадке Ледуа--Вольфа]
$\hat{\Sigma}=(1-\rho)S+\rho\nu I\succ0$ п.н.\ при любом $n\le d$; все тождества выше переносятся на усаженные оценки; при $\rho\to1$ метрика непрерывно деградирует к изотропной, т.е.\ $B1w\to B1$ без ложной геометрии.
\end{lemma}
Численная верификация на закоммиченных кэшах: тождество $\Sigma_T=\Sigma_W+\Sigma_B$ с относительной точностью следа $10^{-9}$--$10^{-8}$; ошибка Шермана--Моррисона 0{,}05--0{,}11 (мера рассинхронизации усадки между $T$- и $W$-путями, качественный вывод не затрагивает).

\secfmt{4. Методология экспериментов}

\textbf{Протокол.} Строгий leakage-free holdout: вся калибровка (средние, ковариации Ледуа--Вольфа, ZCA, SVD-базис, пороги)~--- только на референсной выборке; тест изолирован через разность множеств. Бюджет: $N_{ref}=200$ на класс (суммарно 400) для Wild/ToxicChat; для XSTest (200~unsafe/250~safe)~--- $\lfloor N_{class}/3\rfloor$: $N_{ref}^{mal}=66$, $N_{ref}^{safe}=83$, тест $134+167=301$. Пять фиксированных сидов (0--4); сравнения~--- парный тест ДеЛонга (средняя разница AUC, победы по сидам, $p$-value). Каждая ячейка отчёта пересчитывается автотестом \path|tests/test_report_consistency.py|; всего \textbf{248~автотестов} (консистентность + RSFI + юнит-тесты геометрии/отбеливания/фильтра), все зелёные.

\textbf{Методы.} $A1$ наивный косинус (raw); $A1b$ косинус после $\Sigma_T$-отбеливания; $A2$ RSFI-SVD ($k=20$, raw); $A3$ SVD после отбеливания; $B1$ дискриминантное среднее (raw); $B1b$ дискриминант после $\Sigma_T$-отбеливания; $B1w$ дискриминант после $\Sigma_W$-отбеливания (пул внутриклассовых ковариаций, центрирование общим референсным средним); $C1$ логистическая регрессия (supervised-потолок), $C1b$ её отбелённая версия; 1-классовый и 2-классовый $k$-NN (среднее top-5 косинусов); внешние zero-shot классификаторы ProtectAI DeBERTa-v3-prompt-injection-v2 и unitary/toxic-bert на тех же сплитах.

\textbf{Эмбеддеры.} \texttt{all-mpnet-base-v2} ($d=768$), \texttt{BAAI/bge-base-en-v1.5} ($d=768$), \texttt{BAAI/bge-large-en-v1.5} ($d=1024$), \texttt{Qwen/Qwen3-Embedding-8B} ($d=4096$).

\secfmt{5. Результаты}

\subsecfmt{5.1. Гетерогенный режим: доминирует safe-aware поправка}

На WildChat (табл.~1) $B1$ даёт +8.2~п.п.\ (mpnet), +9.1~п.п.\ (bge-base), +10.5~п.п.\ (bge-large) над наивным косинусом, 5/5~сидов, $p<0{,}0001$; на Qwen3-8B: $B1=0{,}8747$ против $A1=0{,}7982$. Разрыв до логистической регрессии~--- около одного пункта ($0{,}8668$ против $0{,}8766$ на mpnet).

\vspace{4pt}\noindent{\small\bfseries Таблица 1. WildChat: ROC-AUC, $N_{ref}=200$ на класс, 5~сидов, mean $\pm$ std.}\par\vspace{2pt}

\begin{center}\footnotesize\renewcommand{\arraystretch}{0.90}
\begin{tabular}{lrrrr}
\toprule
Модель & $A1$ & $A2$ (SVD) & $\mathbf{B1}$ & $C1$ (LogReg) \\
\midrule
mpnet-base & $0{,}7846\pm0{,}0047$ & $0{,}7875\pm0{,}0074$ & $\mathbf{0{,}8668\pm0{,}0060}$ & $0{,}8766\pm0{,}0071$ \\
bge-base & $0{,}7708\pm0{,}0040$ & $0{,}7854\pm0{,}0114$ & $\mathbf{0{,}8618\pm0{,}0050}$ & $0{,}8728\pm0{,}0038$ \\
bge-large & $0{,}7674\pm0{,}0065$ & $0{,}7839\pm0{,}0110$ & $\mathbf{0{,}8719\pm0{,}0050}$ & $0{,}8804\pm0{,}0028$ \\
Qwen3-8B & $0{,}7982\pm0{,}0102$ & $0{,}8012\pm0{,}0130$ & $\mathbf{0{,}8747\pm0{,}0090}$ & $0{,}8847\pm0{,}0083$ \\
\bottomrule
\end{tabular}
\end{center}

\subsecfmt{5.2. Разделение эффектов: режимы решают всё}

Полная картина (табл.~2, 4~эмбеддера $\times$ 3~датасета): на омонимическом XSTest весь прирост даёт отбеливание ($A1b=0{,}8973\approx B1b=0{,}8970$, $\Delta=-0{,}0003$, $p=0{,}2819$~--- safe-aware незначим); на Wild~--- наоборот, safe-aware в сыром пространстве, а тотальное отбеливание вредит ($0{,}8668\to0{,}8297$ на mpnet); на гомогенном ToxicChat работают оба механизма, причём $B1b$ превосходит отбелённую LogReg ($0{,}9617$ против $0{,}9581$, $p=0{,}0055$).

\vspace{4pt}\noindent{\small\bfseries Таблица 2. ROC-AUC по режимам (среднее 5~сидов; фрагмент: mpnet и Qwen3-8B).}\par\vspace{2pt}

\begin{center}\footnotesize\renewcommand{\arraystretch}{0.90}
\begin{tabular}{llrrrr}
\toprule
Датасет & Модель & $A1$ & $B1$ & $B1b$ ($\Sigma_T$) & $B1w$ ($\Sigma_W$) \\
\midrule
ToxicChat & mpnet & 0{,}9158 & 0{,}9509 & 0{,}9617 & \textbf{0{,}9680} \\
 & Qwen3-8B & 0{,}8082 & 0{,}9628 & 0{,}9679 & \textbf{0{,}9716} \\
Wild & mpnet & 0{,}7846 & \textbf{0{,}8668} & 0{,}8297 & 0{,}8475 \\
 & Qwen3-8B & 0{,}7982 & \textbf{0{,}8747} & 0{,}8224 & 0{,}8391 \\
XSTest & mpnet & 0{,}7618 & 0{,}7851 & 0{,}8970 & \textbf{0{,}8999} \\
 & Qwen3-8B & 0{,}7824 & 0{,}8269 & \textbf{0{,}9381} & 0{,}9370 \\
\bottomrule
\end{tabular}
\end{center}

\subsecfmt{5.3. Внутриклассовое отбеливание}

$B1w\ge B1b$ в 11~из 12~конфигураций; единственное исключение~--- XSTest~$\times$~Qwen3-8B ($0{,}9370$ против $0{,}9381$, $|\Delta|=0{,}0011$ в пределах per-seed std $0{,}0115$). Восстановление потерь на Wild: 48{,}0\% (mpnet), 56{,}4\% (bge-base), 39{,}3\% (bge-large), 31{,}9\% (Qwen3-8B)~--- в среднем \textbf{43{,}9\%}; превосходство на Wild значимо 5/5~сидов у всех эмбеддеров. На ToxicChat $B1w$ достигает 0{,}9680--0{,}9716, на XSTest~--- 0{,}8999--0{,}9370.

\subsecfmt{5.4. Сравнение с $k$-NN и внешними классификаторами}

Направление $B1$ хранит ровно один вектор ($O(d)$, замер: косинусный скор $\approx0{,}003$~мс) против $O(N_{ref}\cdot d)$ у $k$-NN ($\approx0{,}005$~мс): на ToxicChat $B1b$ обходит оба $k$-NN на всех эмбеддерах; на Wild 2-классовый $k$-NN сильнее сырого $B1$ на +0{,}007--0{,}011~AUC ценой 400-векторного индекса; на XSTest оба $k$-NN проигрывают отбелённому дискриминанту ($0{,}84$ против $0{,}90$ на mpnet). Внешние классификаторы на общих сплитах (табл.~3): $B1$ превосходит оба на всех трёх датасетах, включая Wild ($0{,}8668$ против $0{,}8405$ у DeBERTa-v3) при задержке на три порядка ниже (скоринг DeBERTa~--- единицы--десятки мс против $\approx0{,}003$~мс). Оговорка честности: обе внешние модели обучены не на джейлбрейках (prompt-injection и гражданская токсичность) и взяты zero-shot без дообучения; вопрос~--- «конкурентоспособен ли лёгкий фильтр с готовым классификатором», а не «лучший детектор в принципе». Симптом доменной несовместимости: DeBERTa-v3 на XSTest даёт AUC ниже случайного ($0{,}4022$)~--- инвертированный сигнал на омонимии, где отбеливание корректно даёт $0{,}90$.

\vspace{4pt}\noindent{\small\bfseries Таблица 3. Внешние классификаторы: ROC-AUC на общих сплитах, 5~сидов.}\par\vspace{2pt}

\begin{center}\footnotesize\renewcommand{\arraystretch}{0.90}
\begin{tabular}{lrrrr}
\toprule
Датасет & DeBERTa-v3 & toxic-bert & $B1$ & $B1b$ \\
\midrule
ToxicChat & $0{,}5698\pm0{,}0100$ & $0{,}7934\pm0{,}0126$ & $0{,}9509\pm0{,}0048$ & $\mathbf{0{,}9617\pm0{,}0020}$ \\
Wild & $0{,}8405\pm0{,}0036$ & $0{,}7236\pm0{,}0087$ & $\mathbf{0{,}8668\pm0{,}0060}$ & $0{,}8297\pm0{,}0090$ \\
XSTest & $0{,}4022\pm0{,}0123$ & $0{,}6416\pm0{,}0178$ & $0{,}7851\pm0{,}0203$ & $\mathbf{0{,}8970\pm0{,}0132}$ \\
\bottomrule
\end{tabular}
\end{center}

\subsecfmt{5.5. Кросс-домен, операционные точки, E11}

Кросс-доменный перенос (train на одном датасете из пяти, тест на остальных, 4~эмбеддера $\times$ 5~сидов): $B1$~--- лучший метод и in-domain (0{,}957), и в переносе (0{,}783); полный отбеливающий конвейер $B1b$ в переносе не лучше сырого косинуса (0{,}748 против 0{,}754)~--- отбеливание требует доменной рекалибровки. Диагональ in-domain после устранения найденной при аудите утечки (перекрытие test/ref давало артефакт $\approx0{,}999$) честна и независимо совпала с E14. Операционные точки на AdvBench/HarmBench: на HarmBench наивный косинус непригоден (TPR@FPR1\% $\approx0{,}32$, FPR@TPR90 до 0{,}46), а семейство $B1$ даёт TPR@FPR1\% 0{,}86--0{,}996 при FPR@TPR90 $\approx0{,}00$--0{,}03 (напр., bge-large $B1w$: AUC~0{,}9989, TPR@FPR1\%~0{,}979). Валидация RSFI в касательном пространстве (E11): см.\ раздел~3; калибровка порога~--- по валидационному safe-пулу.

\secfmt{6. Границы применимости}

\subsecfmt{6.1. Статическая обфускация (E6b)}

Калибровка~--- только на чистых атаках; тест~--- base64/leetspeak/ROT13/zero-width/homoglyph. Иерархия деградации на Wild при base64 (по трём эмбеддерам): $A1$ $-0{,}516$, $A2$ $-0{,}574$ (семейство разрушено полностью), сырой $B1$ $-0{,}177$~--- самый устойчивый из геометрических, сопоставим с supervised-потолком $C1$ ($-0{,}173$); абсолютный уровень: $B1$~0{,}705 против $0{,}238$ у наивного косинуса (mpnet). Тип обфускации меняет оптимальную предобработку (под leetspeak лучше отбелённые варианты, под base64~--- сырой $B1$); единой конфигурации нет. Zero-width~--- no-op для рассмотренных стеков (максимальное отклонение $1{,}1\cdot10^{-4}$, токенизаторы нормализуют ZWSP). XSTest~--- худший случай: все методы падают до $\le0{,}61$. Вывод: граница относится ко всему эмбеддинг-конвейеру (потолок $C1$ теряет столько же), закрывается L1-декодером (unwrapper) перед эмбеддингом.

\subsecfmt{6.2. Внешние модели под обфускацией (E9b)}

Рост их AUC под трансформациями (до $0{,}90$--1{,}00)~--- OOD-артефакт («аномальные токены $\Rightarrow$ инъекция»), а не робастность; фиксируем как артефакт измерения.

\subsecfmt{6.3. Адаптивный осведомлённый противник (E6c)}

Линейный guardrail $B1(x)=\langle\tilde{e}(x),\mathbf{w}_{disc}\rangle$ имеет коразмерность~1 ($d-1$ степеней свободы у противника). Жадная замена синонимов против $B1$ обходит его при изменении \textbf{0{,}5\%} слов на Wild (ASR@1\%FPR 98{,}3\%), 2{,}9\% на ToxicChat (72{,}0\%), 4{,}5\% на XSTest (81{,}7\%). Атака не переносится на $B1w$ (ASR 49{,}7\% / 7{,}0\% / 2{,}0\%) и на внешние нейросетевые модели (DeBERTa-v3 сохраняет AUC~0{,}85). Причина~--- рассинхронизация градиентов: $\nabla s_{B1}\propto\mathbf{w}_{raw}$, $\nabla s_{B1w}\propto\Sigma_W^{-1}\mathbf{w}_{raw}$, косинус между ними $\approx0{,}58$ (угол $\approx54{,}4^\circ$). Прямая атака на сам $B1w$ его обходит (ASR 90{,}3\% Wild, 25{,}0\% ToxicChat), подтверждая общность ограничения коразмерности~1: $\Sigma_W$-отбеливание защищает от переноса суррогатного градиента, но не от осведомлённого таргетирования.

\secfmt{7. Обсуждение и рекомендации}

Универсальное правило предобработки: при наличии двух классов калибровать отбеливание по пулу внутриклассовой ковариации $\Sigma_W$, а не по объединённой $\Sigma_T$ ($\Sigma_W$ не уступила $\Sigma_T$ ни в одной из 12~конфигураций). Доменно-агностичный бейзлайн по умолчанию~--- сырой $B1$ (лучший и in-domain, и в переносе). Рекомендуемая роль: ультрабыстрый эшелон L1 в defense-in-depth ($\text{L1 Unwrapper}\to\text{L2 Geometric Pre-filter}\to\text{L3 LLM-Judge/Classifier}$), отсекающий 17--33\% атак при FPR~$\le1\%$ без задержки инференса; полная стоимость пайплайна определяется эмбеддером (локальный ONNX~--- 5--10~мс, удалённый API~--- 100--300~мс), сам скоринг~--- микросекунды. Для уровня ВАК~К1/Q1 потребуется сквозное продакшен-тестирование в связке с генеративными судьями (Llama Guard~3, ShieldLM)~--- фиксируем как открытую работу, а не как заявку.

\secfmt{8. Заключение}

\begin{enumerate}\setlength{\itemsep}{0pt}\setlength{\parsep}{0pt}
\item Одноклассовые геометрические фильтры теряют качество из-за игнорирования безопасного класса; дискриминант $\mu_{mal}-\mu_{safe}$ закрывает разрыв (+7.7--10.5~п.п.\ на Wild, $p<0{,}0001$), достигая уровня supervised-потолка при одном векторе и $O(d)$.
\item Доказано разложение $\Sigma_T=\Sigma_W+\Sigma_B$ и формула подавления сигнала; $\Sigma_W$-отбеливание восстанавливает в среднем 43{,}9\% потерь и не уступает $\Sigma_T$ нигде.
\item Лёгкий фильтр превосходит опубликованные трансформерные классификаторы на общих сплитах при задержке на три порядка ниже.
\item Границы зафиксированы честно: base64, адаптивный противник с коразмерностью~1, доменная чувствительность отбеливания~--- с конструктивными закрытиями (unwrapper, рекалибровка, defense-in-depth).
\end{enumerate}

\secfmt{Список литературы / References}
\begingroup\footnotesize

\begin{enumerate}\setlength{\itemsep}{0pt}\setlength{\parsep}{0pt}
\item Ledoit O., Wolf M. A well-conditioned estimator for large-dimensional covariance matrices. \emph{J. Multivariate Analysis}, 2004, 88: 365--411.
\item Su J. et al. Whitening sentence representations for better semantics and faster retrieval. \emph{Proc. ACL}, 2021.
\item Ethayarajh K. How contextual are contextualized word representations? \emph{Proc. EMNLP}, 2019.
\item Absil P.-A. et al. \emph{Optimization Algorithms on Matrix Manifolds}. Princeton Univ. Press, 2008.
\item Boumal N. \emph{An Introduction to Optimization on Smooth Manifolds}. Cambridge Univ. Press, 2023.
\item Lee J. M. \emph{Introduction to Riemannian Manifolds}. Springer, 2018.
\item Arditi A. et al. Refusal in language models is mediated by a single direction. \emph{Proc. NeurIPS}, 2024.
\item Zhao Y. et al. LLMs encode harmfulness and refusal separately. \emph{Proc. NeurIPS}, 2025.
\item Liu C. et al. TrajGuard: streaming hidden-state trajectory detection. \emph{Findings of ACL}, 2026.
\item Zou A. et al. Representation engineering: a top-down approach to AI transparency. \emph{arXiv:2310.01405}, 2023.
\item Inan H. et al. Llama Guard: LLM-based input-output safeguard. \emph{arXiv:2312.06674}, 2023.
\item Reuter D. et al. NeMo Guardrails. \emph{Proc. EMNLP}, 2023.
\item Dong Y. et al. Building guardrails for large language models: a position paper. \emph{Proc. ICML}, 2024.
\item Sekar A. et al. Zero-shot embedding drift detection. \emph{NeurIPS Lock-LLM Workshop}, 2025.
\item Alanova S. et al. Cross-lingual jailbreak detection via semantic codebooks. \emph{arXiv:2604.25716}, 2026.
\item Han P. et al. SafeSwitch: steering unsafe LLM behavior. \emph{Proc. EMNLP}, 2025.
\item He Z. et al. Interpretable LLM guardrails via sparse representation steering. \emph{arXiv:2503.16851}, 2025.
\item Lin Z. et al. ToxicChat: unveiling hidden challenges of toxicity detection. \emph{Findings of EMNLP}, 2023.
\item R\"ottger P. et al. XSTest: a test suite for identifying exaggerated safety behaviours. \emph{Proc. NAACL}, 2024.
\item Zou A. et al. Improving alignment and robustness with circuit breakers. \emph{arXiv:2307.15043}, 2023.
\item Mazeika M. et al. HarmBench. \emph{arXiv:2402.05793}, 2024.
\item Taori R. et al. Stanford Alpaca. GitHub, 2023.
\item Sun X., Xu W. Fast implementation of DeLong's algorithm. \emph{IEEE Signal Proc. Lett.}, 2014.
\item Athalye A. et al. Obfuscated gradients give a false sense of security. \emph{Proc. ICML}, 2018.
\item Reimers N., Gurevych I. Sentence-BERT. \emph{Proc. EMNLP}, 2019.
\item Perez F., Ribeiro I. Ignore this title and hack this review. \emph{arXiv:2211.09527}, 2022.
\end{enumerate}

\endgroup

\end{document}
